\section{Gibbs 不等式与 KL 的方向性}
\label{sec:07-Information-Theory/03-Divergences/02}

上一节用编码代价解释了 KL 散度。那个解释给出一个强烈的暗示：

\[
\KL(P \| Q) \ge 0,
\]

并且等号只在 \(P=Q\) 时成立。暗示归暗示，KL 的求和里每一项并不都非负。当 \(Q(x) > P(x)\) 时，\(\log[P(x)/Q(x)] < 0\)——这是一个负数乘以一个正权重 \(P(x)\)。为什么所有正负项加起来不可能得到负数？这个担保需要一个证明。

\subsection{从 \(\ln u \le u-1\) 出发}

对任意 \(u > 0\)，函数 \(\ln u\) 位于其 \(u=1\) 处的切线下方：

\[
\ln u \le u - 1,
\]

等号仅当 \(u = 1\)。这个不等式可以从 \(\ln u\) 的凹性直接看出：曲线在每一点都不超过自己的切线。取 \(u = Q(x)/P(x)\)，两边乘以 \(P(x)\)（假设 \(P(x) > 0\) 的位置），再对所有 \(x\) 求和：

\[
\begin{aligned}
\KL(P \| Q)
&= \sum_x P(x) \log \frac{P(x)}{Q(x)} \\
&= -\sum_x P(x) \ln \frac{Q(x)}{P(x)} \quad (\text{换成自然对数，只差正常数因子}) \\
&\ge -\sum_x P(x) \left( \frac{Q(x)}{P(x)} - 1 \right) \\
&= \sum_x P(x) - \sum_{x: P(x)>0} Q(x).
\end{aligned}
\]

第一个和等于 1。第二个和：因为 \(Q\) 是在 \(P\) 支持上的概率质量，它不超过 \(Q\) 的总质量 1。所以

\[
\sum_x P(x) - \sum_{x: P(x)>0} Q(x) \ge 1 - 1 = 0.
\]

如果 \(Q\) 还在 \(P\) 支持之外分配了质量，第二个和严格小于 1，结论依然是 \(\ge 0\)。

要取得等号，两步都不能松：首先需要所有 \(P(x) > 0\) 的位置上 \(Q(x)/P(x) = 1\)，即 \(Q(x) = P(x)\)；其次需要 \(Q\) 没有在 \(P\) 支持之外放置任何概率。两项合起来就是 \(P = Q\)（离散情形逐点相等）。

这就是 Gibbs 不等式：

\[
\KL(P \| Q) \ge 0, \qquad \KL(P \| Q) = 0 \iff P = Q.
\]

如果存在 \(P(x) > 0\) 而 \(Q(x) = 0\) 的位置，KL 已经是 \(+\infty\)（上一节的支持条件），结论照样成立。

\subsection{另一条路：直接用 Jensen}

\(-\ln u\) 是凸函数。把 KL 写成期望：

\[
\KL(P \| Q) = \E_P\!\left[-\ln \frac{Q(X)}{P(X)}\right].
\]

Jensen 不等式说凸函数的期望不小于期望处的函数值：

\[
\KL(P \| Q) \ge -\ln \E_P\!\left[\frac{Q(X)}{P(X)}\right].
\]

如果 \(Q\) 的全部质量都在 \(P\) 的支持上，内层期望等于 \(\sum_x Q(x) = 1\)，右边是 \(-\ln 1 = 0\)。如果 \(Q\) 还在支持外分配质量，期望小于 1，右边为正。这个版本更短，但等号条件的讨论不如前一个证明透明——你需要检查 Jensen 取等的条件，还要追踪支持外质量的影响。

两条路各有用途。第一个证明让你看清每一步里支持条件和等号如何联动；第二个证明揭示了 KL 非负不过是 \(-\ln\) 凸性换了一套加权方式写出来。

\subsection{一个不等式，多项事实}

Gibbs 不等式不只是 KL 非负的证明。把不同的分布代入 \(P\) 和 \(Q\) 的位置，立刻得到信息论里的一组基本性质。

\subsubsection{真实模型使交叉熵最小}

从分解式 \(H(P, Q) = H(P) + \KL(P \| Q)\) 和 KL 非负直接得到

\[
H(P, Q) \ge H(P),
\]

等号仅当 \(Q = P\)。用损失函数的语言说：对数损失是一种严格恰当评分规则。如果预测者希望最小化长期期望损失，诚实报告自己的真实概率判断是最优策略——任何系统性偏离都会增加期望损失。

\subsubsection{均匀分布使熵最大}

设字母表大小为 \(m\)，令 \(U(x) = 1/m\) 为均匀分布：

\[
\begin{aligned}
\KL(P \| U)
&= \sum_x P(x) \log \frac{P(x)}{1/m} \\
&= \log m - H(P) \ge 0.
\end{aligned}
\]

所以 \(H(P) \le \log m\)，等号仅当 \(P = U\)。``均匀分布最不确定''不再只是一句图像直觉——它是 KL 非负的一个直接推论。

\subsubsection{互信息非负}

把联合分布和边缘乘积代入 KL：

\[
I(X;Y) = \KL(P_{XY} \| P_X P_Y) \ge 0,
\]

等号恰好对应 \(P_{XY} = P_X P_Y\)，即独立。于是就互信息非负这一条，也统一到了 Gibbs 不等式之下。

\subsection{方向性不是可以用数字弥补的差异}

让我们算一个具体例子，感受两个方向的 KL 在数值上可以差多远。取

\[
P = (0.9, 0.1), \qquad Q = (0.5, 0.5).
\]

以 2 为底：

\[
\begin{aligned}
\KL(P \| Q) &= 0.9 \log_2 \frac{0.9}{0.5} + 0.1 \log_2 \frac{0.1}{0.5} \\
&= 0.9 \log_2 1.8 + 0.1 \log_2 0.2 \approx 0.531 \text{ bit}, \\[6pt]
\KL(Q \| P) &= 0.5 \log_2 \frac{0.5}{0.9} + 0.5 \log_2 \frac{0.5}{0.1} \\
&= 0.5 \log_2 0.556 + 0.5 \log_2 5.0 \approx 0.737 \text{ bit}.
\end{aligned}
\]

都是比较同一对分布，但 \(0.531 \neq 0.737\)。如果有人只写``\(P\) 与 \(Q\) 的 KL 是多少''而不标明方向，信息是不完整的。

\subsection{方向性塑造了两种不同的选择策略}

设真实分布 \(P\) 在空间中有两个相距很远的高概率区域，而你的近似模型族只能放置一个较窄的单峰分布。

\begin{itemize}
\tightlist
\item
  最小化 \(\KL(P \| Q)\)：期望在 \(P\) 下取。\(P\) 常出现而 \(Q\) 没覆盖到的区域会被重罚，所以 \(Q\) 倾向于张开到能覆盖所有重要模式——代价是中间的低概率区域也被迫分配了一些概率质量，可能使 \(Q\) 看起来比 \(P\) 更``胖''。
\item
  最小化 \(\KL(Q \| P)\)：期望在 \(Q\) 下取。\(Q\) 可以避开 \(P\) 的低密度区域——反正 \(Q\) 自己不常去那里，不会受罚。因此 \(Q\) 倾向于钻进某个模式内部，宁可在单一模式上做得很好，也不愿跨过中间的低概率地带去够另一个模式。
\end{itemize}

这常被总结为前向 KL``倾向覆盖模式''，反向 KL``倾向寻找模式''。但这是帮助记忆的典型行为描述，不是无条件定理。实际结果还取决于可选分布族、参数化和优化过程的具体约束。

\subsection{数据处理：后处理不会让分布更可分}

KL 散度也有自己的数据处理不等式。让两个分布 \(P, Q\) 都通过同一个随机通道 \(K\)，记输出分布为 \(PK, QK\)。则

\[
\KL(PK \| QK) \le \KL(P \| Q).
\]

直觉：如果你收到一个样本，它来自 \(P\) 还是 \(Q\)？对已经收到的样本再压缩、加噪、聚合，不会凭空增加辨别来源的证据——你距离原始数据更远了。

形式上这可以由 log-sum 不等式证明：对非负序列 \(a_i, b_i\)，

\[
\sum_i a_i \log \frac{a_i}{b_i} \ge \Bigl(\sum_i a_i\Bigr) \log \frac{\sum_i a_i}{\sum_i b_i}.
\]

把通道的输出概率按输入分组求和即得。这与互信息的数据处理不等式是同一条代数原则的不同面孔——一个约束联合分布内的信息流，一个约束分布间的可分辨性。

举一个随手可验的例子。两颗六面骰子，面概率分布分别是 \(P\) 和 \(Q\)，它们之间的 KL 是某个值 \(d\)。现在你只能观察结果的奇偶性——把六种结果粗分成两类。奇偶分布之间的 KL 不可能超过原来的 \(d\)。合并结果类别只会掩盖差异，不会放大差异。

\subsection{非负不等于距离}

Gibbs 不等式只给了两条：非负，以及相同则零。数学上要成为度量，还需要对称性和三角不等式。KL 两条都不满足。所以以下推理是错误的：

\begin{quote}
KL 非负，相同分布时为零，所以它就是两个分布之间的距离。
\end{quote}

方向性对编码、变分推断和假设检验有明确的操作意义——在那些场景里，本来就存在真实分布与近似分布的不对称关系。但在聚类、可视化或需要无方向相似度的任务里，KL 的方向性可能与任务需求错位。

这就引出一个自然的问题：能不能构造一个量，保留 KL 的信息论解释，同时做到对称、有界，并且在支持不重合时仍然有限？下一节给出一个具体方案。
